\documentclass{report}
\input{preamble.tex}
\usepackage[scr]{rsfso}
\usetikzlibrary{circuits.ee.IEC}

\newcommand{\faketarget}{\oplus\!\!\!\!\odot}
\newcommand{\target}{%
  \begin{tikzpicture}[scale=0.5]
    \fill[black] (0,0) circle (0.1);
    \draw (0,0) circle (0.2);
    \draw (0,0) circle (0.3);
  \end{tikzpicture}%
}

\usepackage[clock]{ifsym}
\usepackage[utopia]{mathdesign}


\title{\Huge{Architecture des ordinateurs}\\{IFT1227}\\{\textbf{Introduction}}}
\author{\huge{Franz Girardin}}
\date{\today}
\lstset{inputencoding=utf8/latin1}

            %%%%%%%%%%%%%%%%%  Sect.                          %%%%%%%%%%%%%%%%%%%%%%%%%%%%%%%%%%%%%%%%%%%%%%%%%%%%%%%%%
\usepackage{helvet}
\titleformat{\chapter}
  {\fontfamily{phv}\bfseries\huge} % format
  {}                % label
  {0pt}             % sep
  {\color{myb}\huge}           % before-code



\titleformat{\section}
  {\normalfont\scshape}{\thesection}{1em}{}
\usepackage{circuitikz}


% Customizing the spacing for the chapter titles
\titlespacing*{\chapter}{0pt}{0pt}{20pt}

% Allow hfill in math environment
\newcommand{\specialcell}[1]{\ifmeasuring@#1\else\omit$\displaystyle#1$\ignorespaces\fi}

% Allow you to do the non implication (implication barred)
\newcommand{\notimplies}{%
  \mathrel{{\ooalign{\hidewidth$\not\phantom{=}$\hidewidth\cr$\implies$}}}}



\DeclareRobustCommand{\looongrightarrow}{%
  \DOTSB\relbar\joinrel\relbar\joinrel\relbar\joinrel\rightarrow
}



% Define the command for the NOT gate
\newcommand{\notgate}{%
    \begin{circuitikz}[scale=0.8, transform shape]
        \draw
        (0,0) node[not port] (mynot) {}
        (mynot.in) -- ++(-0.5,0) node[left] {A}
        (mynot.out) -- ++(0.5,0) node[right] {Y};
    \end{circuitikz}
}


\newcommand{\bufgate}{%
    \begin{circuitikz}[scale=0.8, transform shape]
        \draw
        (0,0) node[buffer port] (mynot) {}
        (mynot.in) -- ++(-0.5,0) node[left] {A}
        (mynot.out) -- ++(0.5,0) node[right] {Y};
    \end{circuitikz}
}



\newcommand{\andgate}{%
    \begin{circuitikz}[scale=0.8, transform shape]
        \draw
        (0,0) node[and port] (myand) {}
        (myand.in 1) -- ++(-0.5,0) node[left] {A}
        (myand.in 2) -- ++(-0.5,0) node[left] {B}
        (myand.out) -- ++(0.5,0) node[right] {Y};
    \end{circuitikz}
}


% Define the command for the OR gate
\newcommand{\orgate}{%
    \begin{circuitikz}[scale=0.8, transform shape]
        \draw
        (0,0) node[or port] (myor) {}
        (myor.in 1) -- ++(-0.5,0) node[left] {A}
        (myor.in 2) -- ++(-0.5,0) node[left] {B}
        (myor.out) -- ++(0.5,0) node[right] {Y};
    \end{circuitikz}
}


% Define the command for the XOR gate
\newcommand{\xorgate}{%
    \begin{circuitikz}[scale=0.8, transform shape]
        \draw
        (0,0) node[xor port] (myxor) {}
        (myxor.in 1) -- ++(-0.5,0) node[left] {A}
        (myxor.in 2) -- ++(-0.5,0) node[left] {B}
        (myxor.out) -- ++(0.5,0) node[right] {Y};
    \end{circuitikz}
}

% Define the command for the NOR gate
\newcommand{\norgate}{%
    \begin{circuitikz}[scale=0.8, transform shape]
        \draw
        (0,0) node[nor port] (mynor) {}
        (mynor.in 1) -- ++(-0.5,0) node[left] {A}
        (mynor.in 2) -- ++(-0.5,0) node[left] {B}
        (mynor.out) -- ++(0.5,0) node[right] {Y};
    \end{circuitikz}
}

% Define the command for the NAND gate
\newcommand{\nandgate}{%
    \begin{circuitikz}[scale=0.8, transform shape]
        \draw
        (0,0) node[nand port] (mynand) {}
        (mynand.in 1) -- ++(-0.5,0) node[left] {A}
        (mynand.in 2) -- ++(-0.5,0) node[left] {B}
        (mynand.out) -- ++(0.5,0) node[right] {Y};
    \end{circuitikz}
}

% Define the command for the XNOR gate
\newcommand{\xnorgate}{%
\begin{circuitikz}[scale=0.8, transform shape]
\draw
(0,0) node[xnor port] (myxnor) {}
(myxnor.in 1) -- ++(-0.5,0) node[left] {A}
(myxnor.in 2) -- ++(-0.5,0) node[left] {B}
(myxnor.out) -- ++(0.5,0) node[right] {Y};
\end{circuitikz}
}


% Define the command for the 3-input NOR gate
\newcommand{\northreegate}{%
    \begin{circuitikz}[scale=0.8, transform shape]
        \draw
        (0,0) node[nor port, number inputs=3] (mynor3) {}
        (mynor3.in 1) -- ++(-0.5,0) node[left] {A}
        (mynor3.in 2) -- ++(-0.5,0) node[left] {B}
        (mynor3.in 3) -- ++(-0.5,0) node[left] {C}
        (mynor3.out) -- ++(0.5,0) node[right] {Y};
    \end{circuitikz}
}


% Define the command for the 3-input AND gate
\newcommand{\andthreegate}{%
\begin{circuitikz}[scale=0.8, transform shape]
    \draw
    (0,0) node[and port, number inputs=3] (myand3) {}
    (myand3.in 1) -- ++(-0.5,0) node[left] {A}
    (myand3.in 2) -- ++(-0.5,0) node[left] {B}
    (myand3.in 3) -- ++(-0.5,0) node[left] {C}
    (myand3.out) -- ++(0.5,0) node[right] {Y};
\end{circuitikz}

}



\begin{document}
\maketitle

\pagebreak

\pagebreak
\begin{multicols*}{3}


    \footnotesize


    \paragraph{Couche logique numérique}  
    Constituées de \textit{portes logiques} construite à partir 
    de \textbf{transisteurs} qui prennent un \textbf{signal}
    \textbf{0} ou \textbf{1} et calcule 
    une \textbf{fonction logique} 
    $\mathbb{ET}$, $\mathbb{OU}$ et $\mathbb{NON}$, etc.

    \paragraph{Porte $\mathbb{NON}$}
    \mbox{}\vspace{1em}\\
    \begin{minipage}{\columnwidth}
        \begin{minipage}[b]{0.5\columnwidth}
            \centering
            \notgate
        \end{minipage}%
        \begin{minipage}[b]{0.5\columnwidth}
            \centering
            \raisebox{0.3\height}{
            \renewcommand{\arraystretch}{1.5}
            \begin{tabular}{c|c}
                A & Y \\
                \hline
                1 & 0 \\
                0 & 1 \\
            \end{tabular}
            }
        \end{minipage}
    \end{minipage}
    \[\entouree{$Y= \overline{A}$} \]


    \paragraph{Porte $\mathbb{BUF}$}
    \mbox{}\vspace{1em}\\
    \begin{minipage}{\columnwidth}
        \begin{minipage}[b]{0.5\columnwidth}
            \centering
            \bufgate
        \end{minipage}%
        \begin{minipage}[b]{0.5\columnwidth}
            \centering
            \raisebox{0.3\height}{
            \renewcommand{\arraystretch}{1.5}
            \begin{tabular}{c|c}
                A & Y \\
                \hline
                0 & 0 \\
                1 & 1 \\
            \end{tabular}
            }
        \end{minipage}
    \end{minipage}
    \[\entouree{$Y= A$} \]



    \paragraph{Porte $\mathbb{ET}$}
    \mbox{}\vspace{1em}\\
    \begin{minipage}{\columnwidth}
        \begin{minipage}[b]{0.5\columnwidth}
            \centering
            \andgate
        \end{minipage}%
        \begin{minipage}[b]{0.5\columnwidth}
            \centering
            \raisebox{-0.3\height}{
            \renewcommand{\arraystretch}{1.5}
            \begin{tabular}{c c | c}
                A & B & Y \\
                \hline
                0 & 0 & 0 \\
                0 & 1 & 1 \\
                1 & 0 & 1 \\
                1 & 1 & 1 
            \end{tabular}
            }
        \end{minipage}
    \end{minipage}
    \[\entouree{$Y = AB $} \]



    \paragraph{Porte $\mathbb{OU}$}
    \mbox{}\vspace{1em}\\
    \begin{minipage}{\columnwidth}
        \begin{minipage}[b]{0.5\columnwidth}
            \centering
            \orgate
        \end{minipage}%
        \begin{minipage}[b]{0.5\columnwidth}
            \centering
            \raisebox{-0.3\height}{
            \renewcommand{\arraystretch}{1.5}
            \begin{tabular}{c c | c}
                A & B & Y \\
                \hline
                0 & 0 & 0 \\
                0 & 1 & 0 \\
                1 & 0 & 0 \\
                1 & 1 & 1 
            \end{tabular}
            }
        \end{minipage}
    \end{minipage}
    \[\entouree{$Y = A + B $} \]



    \paragraph{Porte $\mathbb{XOR}$}
    \mbox{}\vspace{1em}\\
    \begin{minipage}{\columnwidth}
        \begin{minipage}[b]{0.5\columnwidth}
            \centering
            \xorgate
        \end{minipage}%
        \begin{minipage}[b]{0.5\columnwidth}
            \centering
            \raisebox{-0.3\height}{
            \renewcommand{\arraystretch}{1.5}
            \begin{tabular}{c c | c}
                A & B & Y \\
                \hline
                0 & 0 & 0 \\
                0 & 1 & 1 \\
                1 & 0 & 1 \\
                1 & 1 & 0 
            \end{tabular}
            }
        \end{minipage}
    \end{minipage}
    \[\entouree{$Y = A \oplus B$} \]
    \columnbreak



    \paragraph{Porte $\mathbb{NAND}$}
    \mbox{}\vspace{1em}\\
    \begin{minipage}{\columnwidth}
        \begin{minipage}[b]{0.5\columnwidth}
            \centering
            \nandgate
        \end{minipage}%
        \begin{minipage}[b]{0.5\columnwidth}
            \centering
            \raisebox{-0.3\height}{
            \renewcommand{\arraystretch}{1.5}
            \begin{tabular}{c c | c}
                A & B & Y \\
                \hline
                0 & 0 & 1 \\
                0 & 1 & 1 \\
                1 & 0 & 1 \\
                1 & 1 & 0 
            \end{tabular}
            }
        \end{minipage}
    \end{minipage}
    \[\entouree{$Y = \overline{AB}$} \]


    \paragraph{Porte $\mathbb{NOR}$}
    \mbox{}\vspace{1em}\\
    \begin{minipage}{\columnwidth}
        \begin{minipage}[b]{0.5\columnwidth}
            \centering
            \norgate
        \end{minipage}%
        \begin{minipage}[b]{0.5\columnwidth}
            \centering
            \raisebox{-0.3\height}{
            \renewcommand{\arraystretch}{1.5}
            \begin{tabular}{c c | c}
                A & B & Y \\
                \hline
                0 & 0 & 1 \\
                0 & 1 & 0 \\
                1 & 0 & 0 \\
                1 & 1 & 0 
            \end{tabular}
            }
        \end{minipage}
    \end{minipage}
    \[\entouree{$Y = \overline{A + B}$} \]



    \paragraph{Porte $\mathbb{XNOR}$}
    \mbox{}\vspace{1em}\\
    \begin{minipage}{\columnwidth}
        \begin{minipage}[b]{0.5\columnwidth}
            \centering
            \xnorgate
        \end{minipage}%
        \begin{minipage}[b]{0.5\columnwidth}
            \centering
            \raisebox{-0.3\height}{
            \renewcommand{\arraystretch}{1.5}
            \begin{tabular}{c c | c}
                A & B & Y \\
                \hline
                0 & 0 & 1 \\
                0 & 1 & 0 \\
                1 & 0 & 0 \\
                1 & 1 & 1 
            \end{tabular}
            }
        \end{minipage}
    \end{minipage}
    \[\entouree{$Y = \overline{A \oplus B}$} \]



    \paragraph{Porte $\mathbb{NOR}3$}
    \mbox{}\vspace{1em}\\
    \begin{minipage}{\columnwidth}
        \begin{minipage}[b]{0.5\columnwidth}
            \centering
            \northreegate
        \end{minipage}%
        \begin{minipage}[b]{0.5\columnwidth}
            \centering
            \raisebox{-0.6\height}{
            \renewcommand{\arraystretch}{1.5}
            \begin{tabular}{c c c | c}
                A & B & C & Y \\
                \hline
                0 & 0 & 0 & 1 \\
                0 & 0 & 1 & 0 \\
                0 & 1 & 0 & 0 \\
                0 & 1 & 1 & 0 \\
                1 & 0 & 0 & 0 \\
                1 & 0 & 1 & 0 \\
                1 & 1 & 0 & 0 \\
                1 & 1 & 1 & 0 
            \end{tabular}
            }
        \end{minipage}
    \end{minipage}
    \[\entouree{$Y = \overline{A + B + C}$} \]
    \columnbreak



    \paragraph{Porte $\mathbb{AND}3$}
    \mbox{}\vspace{1em}\\
    \begin{minipage}{\columnwidth}
        \begin{minipage}[b]{0.5\columnwidth}
            \centering
            \andthreegate
        \end{minipage}%
        \begin{minipage}[b]{0.5\columnwidth}
            \centering
            \raisebox{-0.6\height}{
            \renewcommand{\arraystretch}{1.5}
            \begin{tabular}{c c c | c}
                A & B & C & Y \\
                \hline
                0 & 0 & 0 & 0 \\
                0 & 0 & 1 & 0 \\
                0 & 1 & 0 & 0 \\
                0 & 1 & 1 & 0 \\
                1 & 0 & 0 & 0 \\
                1 & 0 & 1 & 0 \\
                1 & 1 & 0 & 0 \\
                1 & 1 & 1 & 1 
            \end{tabular}
        }         
    \end{minipage}
    \end{minipage}
    \[\entouree{$Y = ABC$} \]

    \paragraph{Définition de la marge de bruit}
    Tolérance d'un circuit aux \textbf{perturbations} 
    pouvant fausser l'interprétation du signal.  
    \begin{itemize}
        \item[$\rhd$] $ V_{IH}$ : $\min(V | V_{in}\Coloneqq \textbf{1})$  
        \item[$\rhd$]  $V_{IL}$ : $\max(V | V_{in}\Coloneqq \textbf{0})$
        \item[$\blacktriangleright$] Signal reçu \textbf{min} ou \textbf{max} 
            est être interprété comme \textbf{1}  ou \textbf{0}. 
        \item[$\rhd$] $ V_{OH}$ : $\min(V | V_{out}\Coloneqq \textbf{1})$  
        \item[$\rhd$] $V_{OL}$ : $\max(V | V_{out}\Coloneqq \textbf{0})$
        \item[$\blacktriangleright$] Signal \textbf{min} ou \textbf{max} 
            que l'émetteur s'engage à fournir pour être interprété 
            comme \textbf{1}  ou \textbf{0}. 
    \end{itemize}
    \begin{align*}
        NM_H = V_{OH} \; (\textit{ém.}) - V_{IH} \; (src.) \\
        NM_L = V_{IL} (src.) - V_{OL} (\textit{ém.})
    \end{align*}

    \begin{figure}[H]
        \begin{center}
            \includegraphics[width=0.30\textwidth]{MargeBruit.png}
        \end{center}
    \end{figure}

    \paragraph{Transistors}
    \mbox{}\vspace{1em}\\
    Éléments de base des circuits électroniques. Ils sont composés de 
    trois broches; le drain, la source, et la \textbf{grille} qui contrôle 
    les deux autres comme un interrupteur.  La figure suivante représente 
    un transistor hors tension \textit{off}, un transistor hors tension 
    mais polarisé \textit{off} et un transistor sous tension \textit{on}.     
    \begin{center}
    \begin{circuitikz}[scale=0.5]
    % Draw NMOS transistor
    \draw (0,0) node[nigfete] (mos) {}
    (mos.gate) node[left] {g}
    (mos.drain) node[above] {d}
    (mos.source) node[below] {s};

    % Draw NMOS in OFF state
    \draw (3,0) node[nigfete] (mosoff) {}
    (mosoff.gate) node[left] {g}
    (mosoff.drain) node[above] {d}
    (mosoff.source) node[below] {s};
    \draw (mosoff.gate) -- ++(0.5,0) node[circ] {} -- (mosoff.source);

    Draw NMOS in ON state
    \draw (6,0) node[nigfete] (moson) {}
    (moson.gate) node[left] {g}
    (moson.drain) node[above] {d}
    (moson.source) node[below] {s};
    \draw (moson.gate) -- ++(0.5,0) node[circ] {} -- (moson.drain);
    \end{circuitikz}
    \end{center}

    \paragraph{Composition d'un circuit} 
    \mbox{}\\ 
    \textbf{Circuit}$\Coloneqq$ E., S., spec. \textit{fonct}., spec. \textit{temp}.   

    \paragraph{Propriétés d'un circ. combinatoire}
    \mbox{} 
    \vspace{0.1em}
    \begin{itemize}
        \item[$\rhd$]  \textbf{Noeud} $\implies$ \textit{In}  
            ou \textbf{connexion à} un \textit{Out}. 
        \item[$\rhd$] \textbf{Aucun} chemin cyclique.    

    \end{itemize}


    \paragraph{Sommes de produits SOP}
    En considérant les variables \textit{In} d'une ligne de la table de vérité, 
    il faut identifier la \textbf{conjonction} (produit) nécessaire 
    pour engendrer un \textbf{1} logique. Un \textbf{minterm} est une 
    représentation du produit engendrant un 1 logique. 

                \begin{table}[H]
                  \begin{center}
                   \renewcommand{\arraystretch}{1.5}
                    \footnotesize
                        \begin{tabular}{|l|l|l||c|}
                        \arrayrulecolor{blue}\hline
                        \rowcolor{lightBlue}
                        \textcolor{myb}{$A$} & \textcolor{myb}{$B$} 
                                           & \textcolor{myb}{$Y$} 
                                           & \textcolor{myb}{\textit{minterm}}  
                        \\
                        \hline
                        \hline
                        \arrayrulecolor{black}
                        $V$ & $F$ & \cellcolor{myr} $0$ & $\overline{AA}$
                        \\
                        \hline
                        $0$ & $1$ & \cellcolor{myg} $1$ & \textcolor{myg}{$\overline{A}B$}
                        \\
                        \hline 
                        $1$ & $0$ & \cellcolor{myr} $0$ & $A\overline{B}$ 
                        \\ 
                        \hline
                        $1$ & $1$ & \cellcolor{myg} $1$ & \textcolor{myg}{$AB$}
                        \\
                        \hline
                        \end{tabular}
                \end{center}
                \end{table}
                \[Y(A, B) = \overline{A}B + AB \leftrightarrow Y(A,B) = \sum (1, 3) \]
    

    \paragraph{Produits de sommes}
    En considérant les variables \textit{In} d'une ligne de la table de vérité, 
    il faut identifier la \textbf{somme} nécessaire 
    pour engendrer un \textbf{0} logique. Un \textbf{maxterm} est une 
    représentation de la somme engendrant un 0 logique. 


                \begin{table}[H]
                  \begin{center}
                   \renewcommand{\arraystretch}{1.5}
                     \footnotesize
                        \begin{tabular}{|l|l|l||c|}
                        \arrayrulecolor{blue}\hline
                        \rowcolor{lightBlue}
                        \textcolor{myb}{$A$} & \textcolor{myb}{$B$} 
                                           & \textcolor{myb}{$Y$} 
                                           & \textcolor{myb}{\textit{maxterm}}  
                        \\
                        \hline
                        \hline
                        \arrayrulecolor{black}
                        $V$ & $F$ & \cellcolor{myg} $0$ & \textcolor{myg}{$A + B$}
                        \\
                        \hline
                        $0$ & $1$ & \cellcolor{myr} $1$ & $A + \overline{B}$
                        \\
                        \hline 
                        $1$ & $0$ & \cellcolor{myg} $0$ & \textcolor{myg}{$\overline{A} + B$} 
                        \\ 
                        \hline
                        $1$ & $1$ & \cellcolor{myr} $1$ & $\overline{A} + \overline{B}$ 
                        \\
                        \hline
                        \end{tabular}
                \end{center}
                \end{table}
                \[Y(A, B) = (A+B)(A + \overline{B}) = \prod (0, 2) \]



    \begin{table}[H]
      \centering
      \renewcommand{\arraystretch}{1.5}
      \setlength{\arrayrulewidth}{0.4pt}
      \arrayrulecolor{blue}
      \footnotesize
      \begin{tabular}{|l|l|l|}
        \hline
        \rowcolor{lightBlue}
        \textcolor{myb}{Axiome} & \textcolor{myb}{Dual} & \textcolor{myb}{Nom} \\
        \hline
        \hline
        \( B = 0 \) if \( B \neq 1 \) & \( B = 1 \) if \( B \neq 0 \) & Binary field \\
        \rowcolor{lightBlue}
        \( \overline{0} = 1 \) & \( \overline{1} = 0 \) & NOT \\
        \( 0 \cdot 0 = 0 \) & \( 1 + 1 = 1 \) & AND/OR \\
        \rowcolor{lightBlue}
        \( 1 \cdot 1 = 1 \) & \( 0 + 0 = 0 \) & AND/OR \\
        \( 0 \cdot 1 \cdot 0 = 0 \) & \( 1 + 0 = 1 + 1 = 1 \) & AND/OR \\
        \hline
        \end{tabular}
    \end{table}



    \begin{table}[H]
      \centering
      \renewcommand{\arraystretch}{1.5}
      \setlength{\arrayrulewidth}{0.4pt}
      \arrayrulecolor{blue}
      \footnotesize
      \begin{tabular}{|l|l|l|}
        \hline
        \rowcolor{lightBlue}
        \textcolor{myb}{Théorème} & \textcolor{myb}{Dual} & \textcolor{myb}{Nom} \\
        \hline
        \hline
        $B \cdot 1 = B$ & $B + 0 = B$ & Identité \\
        \rowcolor{lightBlue}
        \( B \cdot  0 = 0\) & \( B + 1 = 1 \) & Élément nul \\
        \( B \cdot B = B \) & \( B + B = B \) & Indépotence \\
        \rowcolor{lightBlue}
                            & \( \overline{\overline{B}}  = B\) & Involution \\
        \( B \cdot \overline{B} = 0 \) & \( B + \overline{B} = 1 \) & Complément \\
        \hline
        \end{tabular}
    \end{table}


    \paragraph{De Morgan}
    \mbox{}\vspace{0.5em}
    \[ \neg(A + B) = \neg A + \neg B \]
    \[ \neg(A \cdot B) = \neg A \cdot \neg B \]






    





    \end{multicols*}
\end{document}
